\chapter{Market making sur demande de cotation d'obligations et d'ETF}\label{ch:hf:bond-and-etf-request-for-quote-market-making}

Une obligation d'entreprise dont la dernière transaction remonte à trois semaines s'affiche à l'écran, dans une demande émise par une société de gestion auprès de cinq dealers. L'un d'eux répond en
moins d'une seconde, sans intervention humaine, car son modèle a déjà évalué l'obligation à partir de ses voisines. Dans l'enchère de ce chapitre, cinq dealers dont les
estimations d'une telle obligation s'écartent chacune d'un quart de point~; ils se fixent sur des marges d'environ 43~cents~; le gagnant réalise 0.94 cent par demande à laquelle il
répond et remporte 13~\% d'entre elles. Un dealer qui fixerait sa marge comme si son estimation était exacte remporterait 31~\% des demandes et perdrait de l'argent.

\section{Tarification des obligations qui ne se négocient pas}

La plupart des obligations d'entreprise se négocient rarement (livre~2, chapitre~22)~: une émission peut ne donner lieu qu'à quelques transactions par mois, une demande de cotation pouvant être la seule
formation des prix qu'elle connaît en une semaine. La fair value d'un dealer pour une telle obligation se construit, elle n'est pas observée. Trois sources l'alimentent~:

\begin{itemize}
\item \begin{sloppypar}\emph{ses propres transactions}, issues du système de déclaration des transactions, pondérées selon leur récence et débarrassées des valeurs aberrantes pour former un composite (le
\texttt{firm.rfq.composite} du livre~2), puis déplacées depuis par la courbe de l'émetteur~;\end{sloppypar}
\item \emph{comparables}~: les autres obligations du même émetteur, les obligations de notation, de secteur et d'échéance proches, chacune avec sa sensibilité à la cible~;
\item \emph{proxys liquides}~: la courbe du Trésor, les indices de crédit et les fonds indiciels cotés (ETF) sur obligations, qui bougent chaque seconde et portent la part de la valeur
de l'obligation liée au marché dans son ensemble.
\end{itemize}

\texttt{firm.rfqmm.fair\_value} combine les deux premières. Quelle que soit la méthode, l'estimation comporte une erreur, et, pour une obligation qui n'a pas été négociée depuis des semaines,
cette erreur est grande par rapport à la fourchette qu'un dealer peut facturer. Le modèle de ce chapitre la fixe à 25 cents pour 100 de valeur nominale, soit un
écart-type, pour chaque dealer~; l'erreur de chacun lui est propre.

\section{Cotation automatique sur demande de cotation}

\begin{definition}[Cotation automatique (auto-quoting)]\label{def:hf:bond-and-etf-request-for-quote-market-making:auto}
L'\emph{cotation automatique}\index{cotation automatique} (auto-quoting) est la réponse automatique d'un dealer aux demandes de cotation~: un système qui évalue chaque demande à partir de ses modèles
et de son inventaire, décide si et à quel niveau répondre, et envoie la réponse sans trader, dans les limites fixées par le desk, les demandes
qui sortent de ces limites étant confiées à une personne.
\end{definition}

\begin{omfigure}
\begin{tikzpicture}[font=\footnotesize,>=stealth,box/.style={draw,rounded corners,fill=omThm!8,minimum width=2.2cm,minimum height=0.7cm,align=center}]
\node[box] (c) at (0,0) {client\\(vend 5~M\$)};
\node[box] (p) at (3.7,0) {plateforme};
\foreach \i/\y in {1/1.6,2/0.8,3/0,4/-0.8,5/-1.6}{
  \node[box,minimum width=1.6cm] (d\i) at (6.8,\y) {dealer \i};
  \draw[->] (p.east) -- (d\i.west);
}
\node[box,fill=omDef!10] (m) at (10.3,0) {fair value\\$-$ marge\\(auto-quote)};
\draw[->] (d3.east) -- (m.west);
\draw[->] (c) -- node[above,font=\scriptsize] {demande} (p);
\draw[->,dashed] (d1.north) to[bend right=12] (c.north);
\node[font=\scriptsize,fill=white,inner sep=1pt] at (3.4,2.25) {le meilleur cours acheteur gagne~; le gagnant voit le prix de couverture};
\end{tikzpicture}

\omcaption{Une demande de cotation sur une plateforme dealer-to-client~: le client interroge cinq dealers à la fois~; chacun évalue l'obligation à partir de son propre modèle et
répond dans le délai imparti~; le client frappe le meilleur cours acheteur s'il est acceptable~; la plateforme communique au gagnant le \omterm{def:hf:bond-and-etf-request-for-quote-market-making:cover}{prix de couverture}, soit le deuxième meilleur cours.
Source~: livre~2, chapitre~22~; présent chapitre.}\label{fig:hf:bond-and-etf-request-for-quote-market-making:flow}
\end{omfigure}

Hendershott et Madhavan ont étudié l'innovation au cœur du dispositif (\cref{fig:hf:bond-and-etf-request-for-quote-market-making:flow})~: une technologie électronique
qui permet à un investisseur de consulter simultanément de nombreux dealers obligataires. Ils ont montré que les enchères électroniques périodiques à sens unique constituent une source de liquidité viable et
importante, y compris pour des obligations peu échangées, un compromis entre la recherche bilatérale et un carnet d'ordres continu. O'Hara et Zhou,
en combinant les données propriétaires d'une plateforme avec les déclarations réglementaires de transactions, ont constaté que la négociation électronique sur demande de cotation reste assez modeste et
segmentée, mais qu'elle a des effets étendus sur les coûts de transaction et la qualité d'exécution, tant dans la négociation électronique que téléphonique, et sur le marché
interdealers~: un marché en pleine transition.

\begin{dated}{2026-09}{L'année d'une plateforme}\label{dat:hf:bond-and-etf-request-for-quote-market-making:mktx}
MarketAxess a annoncé un chiffre d'affaires record de 846~millions de dollars pour 2025, une hausse de 48~\% du volume quotidien moyen des transactions de portefeuille, qui atteint un record de 1.4~milliard de dollars,
une hausse de 24~\% du volume quotidien moyen des transactions de blocs, qui atteint un record de 5~milliards de dollars, et une hausse de 33~\% du volume initié par les dealers~; environ 2\,100 sociétés utilisent
sa plateforme.
\end{dated}

\section{Probabilité de gain et malédiction du vainqueur}

\begin{definition}[Prix de couverture (cover price)]\label{def:hf:bond-and-etf-request-for-quote-market-making:cover}
Le \emph{prix de couverture}\index{prix de couverture} (cover price) est la deuxième meilleure cotation d'une demande de cotation, que de nombreuses plateformes communiquent au dealer gagnant après
la transaction~; il lui indique de combien il aurait pu réduire sa cotation tout en gagnant.
\end{definition}

\begin{definition}[Modèle de probabilité de gain (win-probability model)]\label{def:hf:bond-and-etf-request-for-quote-market-making:win}
Un \emph{modèle de probabilité de gain}\index{modèle de probabilité de gain} (win-probability model) estime la probabilité que la cotation d'un dealer remporte une demande de cotation, en fonction
de l'écart entre cette cotation et la fair value du dealer et des variables de la demande (taille, sens, client, obligation, temps écoulé depuis la dernière transaction, nombre de
dealers interrogés), ajusté sur l'historique des gains, des pertes et des \omterm{def:hf:bond-and-etf-request-for-quote-market-making:cover}{prix de couverture} du dealer.
\end{definition}

Le dealer choisit sa marge $m$, c'est-à-dire la distance entre son cours acheteur et sa fair value, de manière à maximiser $P(\text{gain}\mid m)\times\mathrm{E}[\text{bénéfice}\mid
\text{gain},m]$. Le second facteur est là où réside la malédiction du vainqueur~: le dealer dont l'estimation est la plus élevée remporte le plus souvent la demande, si bien que son estimation, sachant
qu'il a gagné, est trop haute. Avec cinq dealers dont les erreurs valent chacune 25 cents, la plus haute estimation est en moyenne surévaluée de 29 cents avant toute marge (l'espérance
du maximum de cinq erreurs normales, \texttt{firm.rfq.expected\_max\_normal} du livre~2)~; avec dix, de 38.5 cents.

L'enchère de ce chapitre (\cref{lst:hf:bond-and-etf-request-for-quote-market-making:sim})~: le client vend au meilleur cours acheteur s'il est au moins égal à son
\omterm{def:hf:inventory-models:reservation}{prix de réservation}, soit la valeur de l'obligation diminuée d'une décote d'urgence de moyenne 30 cents. Les concurrents utilisent une marge commune~; la nôtre est la meilleure réponse~;
l'équilibre est la marge pour laquelle la meilleure réponse à l'ensemble des dealers qui l'utilisent est elle-même. Un modèle qui ignore la malédiction maximise $P(\text{gain})\times m$,
comme si le profit sur une victoire était la marge.

\begin{omfigure}
\begin{tikzpicture}
\begin{axis}[width=11cm,height=5.4cm,axis y line*=right,axis x line=none,xmin=0.5,xmax=10.5,ymin=0,ymax=32,
  ylabel={\small taux de victoire (\%)},tick label style={font=\footnotesize},table/col sep=comma]
\addplot[gray,thick,dotted,mark=triangle*,mark size=1.4pt] table[x=n,y=win]{figdata/hft/22-bond-and-etf-request-for-quote-market-making/dealers.csv};
\label{plot:hf:rfqwin}
\end{axis}
\begin{axis}[width=11cm,height=5.4cm,axis y line*=left,xmin=0.5,xmax=10.5,xtick={1,2,3,5,8,10},xlabel={\small dealers interrogés},
  ylabel={\small cents par demande},ymin=-2,ymax=9,tick label style={font=\footnotesize},grid=major,grid style={gray!25},
  table/col sep=comma,legend style={font=\scriptsize,at={(0.5,-0.3)},anchor=north,/tikz/every even column/.append style={column sep=8pt}},legend columns=3]
\addplot[omThm,thick,mark=*,mark size=1.3pt] table[x=n,y=pnl]{figdata/hft/22-bond-and-etf-request-for-quote-market-making/dealers.csv};
\addlegendentry{profit à la marge d'équilibre}
\addplot[omDef,thick,dashed,mark=square*,mark size=1.3pt] table[x=n,y=naive]{figdata/hft/22-bond-and-etf-request-for-quote-market-making/dealers.csv};
\addlegendentry{en ignorant la malédiction}
\addlegendimage{/pgfplots/refstyle=plot:hf:rfqwin}\addlegendentry{taux de victoire (à droite)}
\draw[gray] (axis cs:0.5,0) -- (axis cs:10.5,0);
\end{axis}
\end{tikzpicture}

\omcaption{Profit espéré d'un dealer par demande traitée, en cents pour 100 de valeur nominale, à la marge d'équilibre symétrique et à celle qu'un
modèle ignorant la malédiction du vainqueur choisirait (les concurrents étant à l'équilibre), ainsi que le taux de victoire d'équilibre (à droite), en fonction du nombre de dealers
interrogés~; erreurs d'estimation de 25 cents, décotes d'urgence de moyenne 30 cents, 40\,000 demandes par point. Données~:
\texttt{hf\_rfq.equilibria}.}\label{fig:hf:bond-and-etf-request-for-quote-market-making:dealers}
\end{omfigure}

\begin{center}
\small
\begin{tabular}{@{}lrrrrrr@{}}
\toprule
nombre de dealers interrogés & 1 & 2 & 3 & 5 & 8 & 10\\
\midrule
marge d'équilibre (cents) & 47.5 & 40.0 & 40.0 & 42.5 & 45.0 & 50.0\\
profit par demande (cents) & 7.96 & 3.46 & 1.96 & 0.94 & 0.41 & 0.45\\
taux de victoire & 27.7~\% & 24.0~\% & 19.1~\% & 13.2~\% & 9.1~\% & 6.9~\%\\
erreur du vainqueur sachant une victoire (cents) & 18.8 & 25.6 & 29.7 & 35.4 & 40.5 & 43.5\\
marge ignorant la malédiction (cents) & 37.5 & 30.0 & 27.5 & 25.0 & 22.5 & 25.0\\
son profit par demande (cents) & 7.50 & 2.99 & 1.25 & $-0.07$ & $-1.13$ & $-0.65$\\
\bottomrule
\end{tabular}
\end{center}

Les résultats (voir la \cref{fig:hf:bond-and-etf-request-for-quote-market-making:dealers})~: la marge d'équilibre ne baisse pas avec le nombre de dealers, elle remonte
à partir de trois, parce que la malédiction croît plus vite que la concurrence ne resserre la fourchette~; le profit par demande passe de 8 cents pour un dealer seul à moins d'un demi-cent
avec huit ou dix dealers. Le dealer qui ignore la malédiction cote 20 à 25 cents plus serré, gagne plus de deux fois plus souvent et perd à partir de cinq dealers.
Même seul, il gagne moins, car les offres gagnantes d'un dealer isolé sont aussi sélectionnées par le \omterm{def:hf:inventory-models:reservation}{prix de réservation} du client.

Le \omterm{def:hf:bond-and-etf-request-for-quote-market-making:win}{modèle de probabilité de gain} est appris sur l'historique. L'ajustement d'une courbe logistique à 20\,000 demandes passées, avec des marges réparties de 0 à 80 cents, face
à cinq concurrents à l'équilibre (voir le \cref{lst:hf:bond-and-etf-request-for-quote-market-making:win}) donne
$P(\text{gain})=1/(1+e^{-(0.665-0.0612\,m)})$~: 12.6~\% à la marge d'équilibre, contre 13.2~\% simulé
(voir la \cref{fig:hf:bond-and-etf-request-for-quote-market-making:winmodel}). Le \omterm{def:hf:bond-and-etf-request-for-quote-market-making:cover}{prix de couverture} alimente le modèle~: chaque victoire révèle de combien le dealer aurait pu
réduire sa marge, chaque défaite ne révèle rien, ce qui biaise les estimations naïves de la distribution des cotations des concurrents~; les modèles qui la décrivent sont ajustés en tenant compte de
cette censure.

\begin{omfigure}
\begin{tikzpicture}
\begin{axis}[width=11cm,height=4.8cm,xlabel={\small marge sous la fair value (cents pour 100 de valeur nominale)},ylabel={\small probabilité de gain},xmin=0,
  xmax=80,ymin=0,ymax=0.65,tick label style={font=\footnotesize},grid=major,grid style={gray!25},table/col sep=comma,
  legend style={font=\scriptsize,at={(0.97,0.97)},anchor=north east}]
\addplot[omThm,only marks,mark=*,mark size=1.6pt] table[x=markup,y=empirical]{figdata/hft/22-bond-and-etf-request-for-quote-market-making/winmodel.csv};
\addplot[omDef,thick] table[x=markup,y=fitted]{figdata/hft/22-bond-and-etf-request-for-quote-market-making/winmodel.csv};
\legend{taux de victoire par classes de 10 cents,modèle logistique ajusté}
\end{axis}
\end{tikzpicture}

\omcaption{Le \omterm{def:hf:bond-and-etf-request-for-quote-market-making:win}{modèle de probabilité de gain}~: taux de victoire de 20\,000 demandes historiques regroupées par marge, et courbe logistique ajustée sur celles-ci par la méthode de
Newton, face à cinq dealers à la marge d'équilibre de 42.5 cents. Données~: \texttt{hf\_rfq.win\_model}.}\label{fig:hf:bond-and-etf-request-for-quote-market-making:winmodel}
\end{omfigure}

\section{Demande de cotation (RFQ) sur ETF et liquidité en blocs}

Les institutions négocient de gros blocs de fonds indiciels cotés (ETF) par demande de cotation (RFQ), hors du carnet d'ordres de la bourse. Le cours acheteur du dealer pour un bloc provient
du panier de création (chapitre~12)~: il peut acheter les actions et les racheter, en recevant le panier qu'il revend. Son cours acheteur est celui du panier diminué
des coûts du rachat~: les frais fixes répartis sur les actions d'une part, le coût de vente du panier et la couverture pendant le règlement-livraison du rachat. Pour un fonds
dont le cours acheteur du panier est de 100.00~\$, pour des parts de 50\,000 actions, des frais de 1\,250~\$, 5 points de base pour vendre le panier et 1 pour le couvrir, le dealer peut coter
à 99.915~\$ par action un bloc de n'importe quelle taille tenant dans son bilan. La malédiction du vainqueur est faible ici, car le panier est coté en continu~;
le risque tient à la liquidité du panier lorsque le bloc est volumineux.

\section{Transactions de portefeuille et négociation all-to-all}

Une transaction de portefeuille (livre~2, chapitre~22) cote une liste de plusieurs centaines d'obligations en un seul paquet, souvent face à un ETF qui en détient la plus grande
partie~: le modèle du dealer cote chaque ligne, et la diversification du paquet réduit la malédiction du vainqueur sur l'ensemble, par rapport à la somme de ses
parties. La croissance des transactions de portefeuille signalée dans l'encadré daté montre que le marché choisit ce type de transaction. La négociation all-to-all (livre~2, chapitre~22) permet aux investisseurs de répondre
aux demandes les uns des autres de façon anonyme~: les non-dealers deviennent des fournisseurs de liquidité, et le cotateur automatique leur fait concurrence.

\section{Fiches de stratégie}

\begin{strategyfile}{Répondeur algorithmique aux RFQ sur obligations d'entreprise}\label{strat:hf:bond-and-etf-request-for-quote-market-making:corp}
\sfield{Qui paie, et pourquoi} Des investisseurs qui doivent négocier des obligations qui s'échangent rarement et paient pour le bilan et l'estimation d'un dealer.
\sfield{Instruments et plateformes} Obligations d'entreprise sur des plateformes dealer-to-client.
\sfield{Signal} Une fair value issue des transactions anciennes de l'obligation, des comparables et des proxies liquides~; le \omterm{def:hf:bond-and-etf-request-for-quote-market-making:win}{modèle de probabilité de gain}.
\sfield{Taille et exécution} Répondre automatiquement dans les limites~; fixer la marge à partir du modèle, en intégrant la malédiction du vainqueur~; orienter les demandes volumineuses ou inhabituelles
vers un trader.
\sfield{Coûts} La malédiction (35 cents d'erreur sachant une victoire face à cinq dealers dans le modèle)~; financement de l'inventaire et couverture.
\sfield{Comment elle meurt} Marges fixées comme si l'estimation était juste (avec cinq dealers, le cotateur naïf gagne 31~\% des demandes et perd de l'argent)~; et
un épisode de stress, lorsque toutes les estimations deviennent obsolètes à la fois, comme lors de la dislocation du marché des obligations d'entreprise de mars 2020 (livre~9, chapitre~19).
\sfield{Horizon, capacité, infrastructure} Quelques secondes pour répondre, de quelques jours à quelques semaines pour le débouclage~; les connexions de la plateforme et les déclarations des transactions.
\sfield{Backtest honnête} Les demandes effectivement reçues par le dealer, les cotations des concurrents (censurées par ce qui a été divulgué), les prix
auxquels l'inventaire pouvait être vendu.
\sfield{Sources} Hendershott et Madhavan (2015)~; O'Hara et Zhou (2021)~; ce chapitre.
\end{strategyfile}

\begin{strategyfile}{Market making par RFQ sur ETF}\label{strat:hf:bond-and-etf-request-for-quote-market-making:etf}
\sfield{Qui paie, et pourquoi} Des institutions qui négocient des blocs de fonds plus importants que la profondeur affichée par la bourse.
\sfield{Instruments et plateformes} Des ETF par demande de cotation (RFQ)~; leurs paniers~; le marché primaire.
\sfield{Signal} La valeur exécutable du panier et les coûts de création ou de rachat.
\sfield{Taille et exécution} Coter à partir du panier, moins (ou plus) les coûts du marché primaire~; couvrir immédiatement dans le panier ou via un contrat à terme.
\sfield{Coûts} \omterm{def:hf:etf-market-making:fee}{Frais de création}, fourchette du panier, couverture.
\sfield{Comment elle meurt} Paniers illiquides en période de stress, où le cours acheteur du panier est lui-même une estimation~: les décotes des fonds obligataires de mars 2020 (livre~9,
chapitre~19) et le gel simulé du chapitre~12.
\sfield{Horizon, capacité, infrastructure} De quelques secondes à une journée~; l'heure limite de création du fonds.
\sfield{Backtest honnête} Prix du panier à l'heure de la demande~; grille des frais~; règlement-livraison.
\sfield{Sources} Chapitre~12~; l'encadré daté.
\end{strategyfile}

\begin{strategyfile}{Streaming sur Treasuries et RFQ}\label{strat:hf:bond-and-etf-request-for-quote-market-making:ust}
\sfield{Qui paie, et pourquoi} Des investisseurs qui négocient des Treasuries sur des plateformes dealer-to-client et préfèrent un prix ferme d'un dealer dont l'identité est connue.
\sfield{Instruments et plateformes} Treasuries on-the-run et off-the-run~; streams et demandes sur les plateformes dealer-to-client.
\sfield{Signal} Les carnets d'ordres interdealer (livre~2, chapitre~4) pour les on-the-run~; une courbe ajustée pour les off-the-run.
\sfield{Taille et exécution} Streamer serré sur les on-the-run, couvrir sur le marché interdealer ou sur des contrats à terme~; répondre aux demandes sur les off-the-run à partir de la courbe.
\sfield{Coûts} Fourchettes de couverture~; erreur d'ajustement de la courbe sur les off-the-run.
\sfield{Comment elle meurt} Des concurrents disposant d'un accès interdealer plus rapide, et des journées où le carnet interdealer s'amincit~; le 15 octobre 2014, le rendement à dix ans
a baissé de 16 points de base entre 9\,h\,33 et 9\,h\,39, puis en a presque effacé la totalité vers 9\,h\,45, sans déclencheur apparent et avec une profondeur étonnamment faible
(le rapport officiel conjoint du personnel de 2015).
\sfield{Horizon, capacité, infrastructure} De la milliseconde à quelques heures.
\sfield{Backtest honnête} Prix interdealer à l'heure de chaque demande, et non au cours de clôture de la journée.
\sfield{Sources} Livre~2, chapitre~4~; Joint Staff Report, The U.S. Treasury market on October 15, 2014 (2015).
\end{strategyfile}

\begin{strategyfile}{Fourniture anonyme de liquidité en all-to-all}\label{strat:hf:bond-and-etf-request-for-quote-market-making:a2a}
\sfield{Qui paie, et pourquoi} Des investisseurs dont les demandes atteignent anonymement des non-dealers.
\sfield{Instruments et plateformes} Obligations d'entreprise sur des protocoles all-to-all.
\sfield{Signal} La même fair value, sans identité de client à prendre en compte.
\sfield{Taille et exécution} Répondre là où l'estimation du modèle est la meilleure (obligations ayant des transactions récentes, comparables proches)~; taille réduite.
\sfield{Coûts} Une malédiction plus aiguë qu'avec des clients connus, puisque le flux ne peut pas être segmenté.
\sfield{Comment elle meurt} Sélection adverse due à des demandeurs anonymes mieux informés.
\sfield{Horizon, capacité, infrastructure} De quelques secondes à quelques jours.
\sfield{Backtest honnête} Mouvements ultérieurs des prix des demandes anonymes, par type de demandeur, lorsque la plateforme les communique.
\sfield{Sources} O'Hara et Zhou (2021)~; livre~2, chapitre~22.
\end{strategyfile}

\section{Tutoriel~: cinq dealers et une obligation que personne n'a négociée}\label{tut:hf:bond-and-etf-request-for-quote-market-making:tut}

\begin{tutorial}
\textbf{Objectif.} Trouver la marge d'équilibre d'un auto-quoter face à des concurrents affectés de la même erreur d'estimation, mesurer la malédiction du vainqueur et ajuster
un \omterm{def:hf:bond-and-etf-request-for-quote-market-making:win}{modèle de probabilité de gain}. \textbf{État final~:} le tableau et les trois figures.
\begin{tutsteps}
\item \textbf{L'enchère}~: l'estimation de chaque dealer avec sa propre erreur, le \omterm{def:hf:inventory-models:reservation}{prix de réservation} du client, le meilleur cours acheteur.
\omcode{code/firm/rfqmm/firm_rfqmm.py}{50}{63}{Estimations, marges, choix du client et profit du vainqueur par rapport à la valeur vraie.}\label{lst:hf:bond-and-etf-request-for-quote-market-making:sim}
\item \textbf{L'équilibre} par meilleures réponses sur une grille de marges (\texttt{firm.rfqmm.equilibrium}).
\item \textbf{Le modèle de gain} ajusté sur un historique randomisé.
\omcode{code/firm/rfqmm/firm_rfqmm.py}{97}{108}{Probabilité de gain logistique en fonction de la marge, par la méthode de Newton.}\label{lst:hf:bond-and-etf-request-for-quote-market-making:win}
\item \textbf{Le cours acheteur du bloc d'ETF} à partir du panier de création (\texttt{firm.rfqmm.etf\_rfq\_bid}).
\end{tutsteps}
\textbf{Pour aller plus loin.} Donner aux dealers des erreurs différentes (un meilleur modèle gagne plus souvent et rapporte davantage)~; ajouter l'inventaire à la marge (axes)~;
estimer la distribution des concurrents à partir des \omterm{def:hf:bond-and-etf-request-for-quote-market-making:cover}{prix de couverture}, avec censure.
\end{tutorial}

\section{Construction~: le teneur de marché RFQ}\label{bld:hf:bond-and-etf-request-for-quote-market-making:rfqmm}

\begin{build}
\bfield{Rôle} Valoriser des obligations peu liquides, répondre aux demandes de cotation à la marge qui maximise le profit espéré net de la malédiction du vainqueur, et valoriser les
blocs d'ETF à partir de leurs paniers.
\bfield{Interface} \texttt{fair\_value(prints, ages\_days, comps, comp\_betas, curve\_move, half\_life)},
\texttt{simulate(n\_dealers, our\_markup, n, seed, sigma\_est, comp\_markup, urgency)}, \texttt{optimise}, \texttt{equilibrium}, \texttt{naive\_markup}, \texttt{fit\_win\_model},
\texttt{etf\_rfq\_bid}. Construit sur \texttt{firm.rfq} (livre~2).
\bfield{Règles} Centimes pour 100 de valeur nominale~; le client vend~; profit par rapport à la valeur vraie de l'obligation.
\bfield{Tests d'acceptation} \texttt{code/firm/rfqmm/tests/}~: la fair value fondée sur le composite, dans l'intervalle des comparables~; le profit égale la valeur moins le cours acheteur sur
les transactions remportées~; l'erreur moyenne du vainqueur est positive~; des marges plus élevées font gagner moins~; avec cinq dealers, profit et taux de gain sont inférieurs à ceux d'un seul dealer~; la marge naïve
est plus basse et rapporte moins~; l'ajustement logistique retrouve les coefficients connus~; le cours acheteur de l'ETF calculé à la main.
\bfield{Pour aller plus loin} Dealers hétérogènes~; inventaire et axes~; estimation censurée à partir des \omterm{def:hf:bond-and-etf-request-for-quote-market-making:cover}{prix de couverture}.
\end{build}

\begin{omsources}
\item T.~Hendershott, A.~Madhavan, Click or call~? Auction versus search in the over-the-counter market, \emph{Journal of Finance} 70(1), 2015,
419--447.
\item M.~O'Hara, X.~A.~Zhou, The electronic evolution of corporate bond dealers, \emph{Journal of Financial Economics} 140(2), 2021, 368--390.
\item MarketAxess Holdings, fourth quarter and full year 2025 results, Form 8-K exhibit 99.1, 6 février 2026.
\item US Department of the Treasury, Federal Reserve, SEC and CFTC staff, Joint Staff Report~: The U.S. Treasury market on October 15, 2014, juillet 2015.
\end{omsources}

\section{Exercices}

\begin{exercise}[$\star$]\label{exo:hf:bond-and-etf-request-for-quote-market-making:1}
Le panier d'un fonds a un cours acheteur de 100.00~\$~; chaque unité compte 50\,000 actions~; les frais sont de 1\,250~\$~; vendre le panier coûte 5 points de base et la couverture 1. Quel est le
cours acheteur du bloc par action~?
\end{exercise}

\begin{exercise}[$\star$]\label{exo:hf:bond-and-etf-request-for-quote-market-making:2}
Avec des erreurs d'estimation de 25 centimes et cinq dealers, de combien l'estimation la plus haute est-elle trop élevée en moyenne~?
\end{exercise}

\begin{exercise}[$\star$]\label{exo:hf:bond-and-etf-request-for-quote-market-making:3}
Le modèle ajusté est $P=1/(1+e^{-(0.665-0.0612m)})$. Quelle est la probabilité de gagner avec une marge de 42.5 centimes~? Et à 20~?
\end{exercise}

\begin{exercise}[$\star\star$]\label{exo:hf:bond-and-etf-request-for-quote-market-making:4}
Pourquoi la marge d'équilibre augmente-t-elle quand on passe de cinq à dix dealers~?
\end{exercise}

\begin{exercise}[$\star\star$]\label{exo:hf:bond-and-etf-request-for-quote-market-making:5}
Pourquoi le fait de ne révéler le \omterm{def:hf:bond-and-etf-request-for-quote-market-making:cover}{prix de couverture} qu'au vainqueur biaise-t-il un \omterm{def:hf:bond-and-etf-request-for-quote-market-making:win}{modèle de probabilité de gain} ajusté naïvement~?
\end{exercise}

\begin{exercise}[$\star\star$]\label{exo:hf:bond-and-etf-request-for-quote-market-making:6}
Pourquoi la malédiction du vainqueur est-elle plus faible dans une transaction de portefeuille que dans la somme de ses lignes~?
\end{exercise}

\begin{exercise}[$\star\star\star$]\label{exo:hf:bond-and-etf-request-for-quote-market-making:7}
\emph{Programmation.} Avec \texttt{firm.rfqmm.optimise}, trouver notre meilleure marge face à cinq concurrents à 42.5 centimes lorsque l'erreur de notre estimation est de 15
centimes et celle des concurrents 25. Que gagnons-nous, et combien de demandes remportons-nous~?
\end{exercise}

\begin{exercise}[$\star\star\star$]\label{exo:hf:bond-and-etf-request-for-quote-market-making:8}
\emph{Trouver la faille.} «~Notre auto-quoter remporte 30~\% des demandes avec une marge de 25 centimes, et notre marge moyenne sur les gains est de 25 centimes~: il gagne 25 centimes
par demande remportée.~»
\end{exercise}

\section{Problème~: cinq dealers et une obligation que personne n'a négociée}

\begin{problem}[{Problème du week-end --- cinq dealers et une obligation que personne n'a négociée}]\label{pb:hf:bond-and-etf-request-for-quote-market-making:1}
Un dealer automatise ses réponses aux demandes de cotation (RFQ) portant sur des obligations d'entreprise peu liquides.

\medskip\noindent\textbf{Partie I --- Prix.}
\begin{enumerate}
\item Comment construit-on une fair value pour une obligation qui n'a pas été négociée depuis des semaines~?
\item Définir la \omterm{def:hf:bond-and-etf-request-for-quote-market-making:auto}{cotation automatique}.
\item Que mettent en évidence Hendershott et Madhavan, ainsi que O'Hara et Zhou~?
\item Résumer l'encadré daté.
\end{enumerate}

\medskip\noindent\textbf{Partie II --- L'enchère.}
\begin{enumerate}[resume]
\item Définir le \omterm{def:hf:bond-and-etf-request-for-quote-market-making:cover}{prix de couverture} (cover price) et un \omterm{def:hf:bond-and-etf-request-for-quote-market-making:win}{modèle de probabilité de gain}.
\item Écrire l'objectif du dealer et expliquer la malédiction du vainqueur qui y figure.
\item Quelle est l'ampleur de la malédiction pour cinq et dix dealers avant toute marge~?
\item Comment trouve-t-on l'équilibre~?
\end{enumerate}

\medskip\noindent\textbf{Partie III --- Mesures.}
\begin{enumerate}[resume]
\item Donner la marge d'équilibre, le profit et le taux de gain pour 1, 5 et 10 dealers.
\item Quel est l'effet de ne pas tenir compte de la malédiction~?
\item Dans quelle mesure le modèle logistique est-il bien ajusté~?
\item Valoriser le bloc d'ETF de l'exemple.
\end{enumerate}

\medskip\noindent\textbf{Partie IV --- Le verdict.}
\begin{enumerate}[resume]
\item Énoncer le \emph{résultat nommé}~: le profit par demande et le taux de gain de l'auto-quoter à sa marge optimale, et la manière dont les deux varient avec le nombre de
dealers consultés.
\item Combien de dealers un client doit-il consulter~?
\item Que procure une meilleure estimation que celle des concurrents~?
\item Pourquoi les transactions de portefeuille se développent-elles~?
\item Qu'est-ce qui change dans la négociation all-to-all~?
\item Quelle fiche de stratégie est la plus exposée à la malédiction~?
\item Comment estimer les cotations des concurrents à partir des \omterm{def:hf:bond-and-etf-request-for-quote-market-making:cover}{prix de couverture}~?
\item En une phrase~: que vend un auto-quoter d'obligations~?
\end{enumerate}
\end{problem}

\section{Questions d'entretien}

\begin{interviewq}[$\star$ \iqroles{trader}]\label{iq:hf:bond-and-etf-request-for-quote-market-making:1}
Vous avez remporté une demande à 99.50 et le \omterm{def:hf:bond-and-etf-request-for-quote-market-making:cover}{prix de couverture} était de 99.20. Qu'est-ce que cela vous apprend~?
\end{interviewq}

\begin{interviewq}[$\star\star$ \iqroles{chercheur}]\label{iq:hf:bond-and-etf-request-for-quote-market-making:2}
Comment valoriseriez-vous une obligation qui n'a pas été négociée depuis un mois~?
\end{interviewq}

\begin{interviewq}[$\star\star$ \iqroles{développeur}]\label{iq:hf:bond-and-etf-request-for-quote-market-making:3}
Concevez un auto-quoter qui répond à 50\,000 demandes par jour en 500 millisecondes chacune. Où se situent les risques~?
\end{interviewq}

\begin{interviewq}[$\star\star$ \iqroles{risque}]\label{iq:hf:bond-and-etf-request-for-quote-market-making:4}
Le taux de gain de l'auto-quoter a doublé le mois dernier. Que vérifiez-vous~?
\end{interviewq}

\begin{interviewq}[$\star\star$ \iqroles{trader}]\label{iq:hf:bond-and-etf-request-for-quote-market-making:5}
Un client vous consulte, vous et dix-neuf autres dealers. Comment cotez-vous~?
\end{interviewq}

\begin{interviewq}[$\star\star\star$ \iqroles{chercheur}]\label{iq:hf:bond-and-etf-request-for-quote-market-making:6}
Avec $n$ dealers dont les estimations ont des erreurs normales indépendantes d'écart type $\sigma$, dériver l'erreur espérée de l'estimation la plus haute pour
$n=2$.
\end{interviewq}
